\section{Efficiency considerations}

In a replicated system, the operating system directly locates objects. In the event of a failure, requests are redirected to another server, but in the usual case objects are accessed directly. 

With erasure coding, fragments of objects must be cataloged. Since the bits necessary to reconstruct an object exist on multiple disks, a system must keep track of where each of these fragments are located. When an object is requested from an $m+n$ system, a server looks up the location of at least $m$ fragments. (The details of this process are implementation dependent. A server could, for example, randomly choose $m$ fragments or try to determine the $m$ closest fragments.) If the requests succeed, these $m$ fragments are transferred to a server to be re-assembled into the requested object. Since assembly cannot begin until the last fragment is available, the process would be slowed down if one of the fragments were coming from a location further away than the others.

